\documentclass[12pt,a4paper]{article}
\usepackage[margin=25mm, headheight=15pt, headsep=10pt]{geometry}
\usepackage{fontspec}
\usepackage[table]{xcolor} % Note: table option is required for rowcolors
\usepackage{titlesec}
\usepackage{tocloft}
\usepackage{fancyhdr}
\usepackage{booktabs}
\usepackage{tcolorbox}
\tcbuselibrary{skins, breakable}
\usepackage[colorlinks=true, linkcolor=emerald, urlcolor=emerald, bookmarks=true]{hyperref}

\definecolor{emerald}{RGB}{0,102,51}
\definecolor{darkred}{RGB}{139,0,0}
\definecolor{lightgray}{RGB}{245,245,245}

\usepackage[nil]{babel}
\babelprovide[import, main]{bengali}
\babelprovide[import]{arabic}
\babelfont{rm}[Renderer=HarfBuzz,Scale=1.0]{Noto Serif Bengali}
\babelfont{sf}[Renderer=HarfBuzz]{Noto Sans Bengali}
\babelfont{tt}[Renderer=HarfBuzz]{Noto Sans Bengali}
\babelfont[arabic]{rm}[Renderer=HarfBuzz,Scale=1.1]{Noto Naskh Arabic}

% Section styling
\titleformat{\section}{\Large\bfseries\color{emerald}}{\thesection.}{0.5em}{}
\titleformat{\subsection}{\large\bfseries\color{emerald!85!black}}{\thesubsection.}{0.5em}{}
\titleformat{\subsubsection}{\normalsize\bfseries\color{emerald!70!black}}{\thesubsubsection.}{0.5em}{}
\titlespacing*{\section}{0pt}{18pt}{8pt}

% Fancy Header/Footer setup
\pagestyle{fancy}
\fancyhf{} % clear all header and footer fields
\fancyhead[L]{\color{emerald}\small\textbf{\nouppercase{\leftmark}}}
\fancyfoot[L]{\scriptsize\color{black!50}{\lat{AI}}-সংকলিত, আলেম-যাচাইকৃত নয়}
\fancyfoot[C]{\color{emerald!70!black}\thepage}
\renewcommand{\headrulewidth}{0.4pt}
\renewcommand{\headrule}{\hbox to\headwidth{\color{emerald}\leaders\hrule height \headrulewidth\hfill}}
\renewcommand{\footrulewidth}{0pt}

% Custom boxes
% 1. Ayat/Hadith Box
\newtcolorbox{ayatbox}[1][]{
    enhanced, breakable, 
    colback=emerald!5, 
    colframe=emerald, 
    boxrule=0.5pt, 
    arc=3pt, 
    left=10pt, right=10pt, top=10pt, bottom=10pt,
    #1
}

% 2. Quote/Translation Box
\newtcolorbox{quotebox}[1][]{
    enhanced, breakable,
    frame hidden,
    colback=lightgray,
    borderline west={3pt}{0pt}{emerald},
    left=15pt, right=10pt, top=10pt, bottom=10pt,
    sharp corners,
    #1
}

% 3. AI-generated notice box
\newtcolorbox{warnbox}[1][]{
    enhanced, breakable,
    colback=red!4,
    colframe=darkred,
    boxrule=0.8pt,
    arc=3pt,
    left=10pt, right=10pt, top=10pt, bottom=10pt,
    #1
}

% Commands
\newcommand{\arab}[1]{\foreignlanguage{arabic}{#1}}
\newcommand{\kib}[1]{\textcolor{emerald}{\textbf{#1}}}

% Latin (English) fallback: Noto Serif Bengali has NO Latin glyphs
\newfontfamily\latfont[Renderer=HarfBuzz]{Noto Serif}
\newcommand{\lat}[1]{{\latfont #1}}

\setlength{\parskip}{5pt}
\setlength{\parindent}{0pt}

% Fix for missing bullet in Noto Serif Bengali
\renewcommand{\labelitemi}{$\bullet$}



\begin{document}
\begin{center}{\Large\bfseries\color{emerald} 11-উমার-ছেঁড়া-জামা-ও-বিনয়}\end{center}
\vspace{1em}
\section*{ঘটনা ১১: উমার (রাঃ) — ছেঁড়া জামা ও ক্ষমতার বিনয়}

\begin{warnbox}
 \textbf{গুরুত্বপূর্ণ নোটিশ (\lat{Important} \lat{Notice}):} এই কন্টেন্টটি \lat{AI} (কৃত্রিম বুদ্ধিমত্তা) দিয়ে প্রস্তুত ও সংকলিত। \lat{AI} কঠোর সূত্র-মান নীতি মেনে শুধু নির্ভরযোগ্য উৎস থেকে তথ্য সংগ্রহ করেছে — কেবল সহীহ/হাসান হাদিস ও সহীহ সনদের আসার রাখা হয়েছে, দুর্বল সনদ বাদ দেওয়া হয়েছে। তবে এটি কোনো আলেম বা মুহাদ্দিস কর্তৃক যাচাইকৃত নয়।
\end{warnbox}


\medskip\hrule\medskip

\subsection*{১. সূত্র কার্ড}

\begin{table}[h]\centering\small
\begin{tabular}{ll}
বিষয় & বিবরণ \\
\hline
\textbf{বর্ণনাকারী} & মুহাম্মদ ইবনে সিরিন (রহ.) থেকে \\
\textbf{গ্রন্থ} & মুসান্নাফ ইবনে আবি শায়বাহ, হাদিস ৩৭১৯১ \\
\textbf{অধ্যায়} & মুসান্নাফ — কিতাবুজ-যুহদ (বিনয় ও ত্যাগ) \\
\textbf{সনদের মান} & \textbf{সহীহ সনদের আসার} — ইবনে সিরিনের সূত্রে মুরসাল/কুতুবী \\
\textbf{মূল বিষয়} & খলিফার জামায় তিনটি ছেঁড়া দাগ — ক্ষমতা, ধন ও প্রদর্শন থেকে দূরে থাকা \\
\hline
\end{tabular}
\end{table}

\medskip\hrule\medskip

\subsection*{২. পটভূমি (কে, কখন, কোথায়, কেন)}

\textbf{উমার ইবনুল খাত্তাব (রাঃ)} ইসলামের ইতিহাসে বিনয়ী শাসকের চরম \lat{example} (উদাহরণ) — তার \lat{humility} (বিনয়) ছিল অতুলনীয়। তিনি খলিফা হওয়ার পরও সাধারণ মানুষের মতো \lat{clothing} (পোশাক) পরতেন, বাজারে যেতেন, নিজের \lat{work} (কাজ) নিজেই করতেন।

মদিনার জীবনে একবার উমারের কাছে মানুষ এসে পরামর্শ/শিক্ষা (\lat{guidance}) গ্রহণ করছিল। তার পরনে ছিল একটি সাধারণ \lat{garment} (পোশাক) — যার বিভিন্ন জায়গায় \textbf{তিনটি \lat{torn} (ছেঁড়া) দাগ} ছিল। মুহাম্মদ ইবনে সিরিন — বিশিষ্ট তাবেয়ি — এই \lat{scene} (দৃশ্য) লক্ষ্য করেছেন এবং এটি বর্ণনা করে গেছেন।

\medskip\hrule\medskip

\subsection*{৩. পূর্ণ ঘটনার বর্ণনা}

মুহাম্মদ ইবনে সিরিন (রহ.) বলেন:

\begin{quote}
\textbf{\lat{“}উমারের (রাঃ) পরনে এমন একটি জামা ছিল, যাতে তিনটি ছেঁড়া দাগ (রুক'আ) ছিল।} \\
\textbf{একটি দাগ তাঁর কাঁধের (কাকুল/কলার) এলাকায়,} \\
\textbf{অন্যটি তাঁর পাঁজরের (পেটের পাশের) এলাকায়,} \\
\textbf{আর অন্যটি তাঁর হাতের কনুইয়ের (মারফিক) এলাকায়।\lat{”}}
\end{quote}

এ ছাড়া তিনি তার পরিধেয় পোশাককে নিয়ে কখনো বড়াই করতেন না, মানুষের সামনে রাজকীয় ভাব দেখাতেন না। ক্ষমতার সর্বোচ্চ শিখরে বসেও উমার একটি ছেঁড়া জামাতেই সন্তুষ্ট থাকতেন — তার জীবন ছিল প্রদর্শনবর্জিত, দুনিয়ামুখী নয়, আখিরাতমুখী।

\medskip\hrule\medskip

\subsection*{৪. ধাপে ধাপে বিশ্লেষণ}

\textbf{ধাপ ১ — দৃশ্য:} উমারের পরনে তিনটি \lat{torn} দাগের জামা। কাঁধে, পাঁজরে ও কনুইতে — অর্থাৎ জামাটি পুরনো ও ব্যবহৃত, তবুও তার মনে কোনো \lat{shame} (লজ্জা) ছিল না।

\textbf{ধাপ ২ — বাস্তবতা:} বিশ্বের সবচেয়ে \lat{powerful} (শক্তিশালী) সাম্রাজ্যের \lat{ruler} (শাসক) — রোম ও ফারসির বিজেতা — কিন্তু তাঁর পোশাক একজন সাধারণ শ্রমিকের মতো।

\textbf{ধাপ ৩ — খলিফার সচেতন পছন্দ:} উমার ধনী হওয়ার \lat{opportunity} (সুযোগ) পেয়েছিলেন; বায়তুল মাল (রাষ্ট্রীয় কোষাগার) তার সামনে খোলা ছিল। তবু তিনি সাধারণ পোশাকেই সন্তুষ্ট ছিলেন — এটা ছিল \lat{deliberate} \lat{choice} (সচেতন সিদ্ধান্ত), দারিদ্র্য নয়।

\textbf{ধাপ ৪ — মুহাম্মদ ইবনে সিরিনের সাক্ষ্য:} তাবেয়িদের বিশিষ্ট ইমাম এই দৃশ্য স্মরণ করেছেন — প্রমাণ করে উমারের \lat{humility} সাহাবা ও তাবেয়িনদের মনে গভীর \lat{impression} (ছাপ) ফেলেছে।

\textbf{ধাপ ৫ — আখিরাতের দিকে তাকিয়ে জীবন:} উমারের পোশাকের এই \lat{simplicity} (সরলতা) — দুনিয়ার \lat{temptation} (মোহ) ও প্রদর্শনের \lat{pride} (অহংকার) থেকে মুক্ত জীবনের প্রতীক।

\medskip\hrule\medskip

\subsection*{৫. শব্দের ব্যাখ্যা}

\begin{table}[h]\centering\small
\begin{tabular}{ll}
শব্দ & অর্থ \\
\hline
\textbf{রুক'আ (\arab{رقعة})} & ছেঁড়া জায়গায় লাগানো টুকরা/দাগ \\
\textbf{কাকুল (\arab{كفل}/\arab{كاهل})} & কাঁধের কাছের অংশ \\
\textbf{মারফিক (\arab{مرفق})} & কনুই \\
\textbf{যুবদ (\arab{زهد})} & দুনিয়া থেকে বিমুখতা — অল্পে তুষ্টি \\
\textbf{তাওয়াদু' (\arab{تواضع})} & বিনয় \\
\hline
\end{tabular}
\end{table}

\textbf{উমারের পোশাক — প্রদর্শনের প্রত্যাখ্যান:} উমার শুধু বিনয়ী ছিলেন না; তিনি নিজেকে রাজকীয় পোশাকের অহংকার থেকে মুক্ত রেখেছিলেন। পোশাক যতই সামান্য হোক, তাতে কোনো লজ্জা নেই — কারণ মর্যাদা তাকওয়ায়, পোশাকে নয়।

\medskip\hrule\medskip

\subsection*{৬. সংশ্লিষ্ট হাদিস ও আয়াত}

\textbf{হাদিস (মুসলিম ৯১):} নবীজি বলেছেন: \textbf{\lat{“}যার হৃদয়ে সরিষার দানার পরিমাণ কিবর থাকবে, সে জান্নাতে প্রবেশ করবে না।\lat{”}}

\textbf{হাদিস (মুসলিম ৯১বি):} \textbf{\lat{“}আল্লাহ সুন্দর (পোশাক-পরিচ্ছদ) গ্রহণযোগ্য বলেছেন, তবে প্রদর্শন (অহংকার) নয়।\lat{”}} — ইবনে মাসউদের বর্ণনায়: পোশাক সুন্দর হতে পারে, কিন্তু তাতে গর্ব করা যাবে না।

\textbf{হাদিস (তিরমিজি ২৩৬৩):} নবীজি বলেছেন: \textbf{\lat{“}যে ব্যক্তি দুনিয়ায় বিনয়ী হয়, আল্লাহ তাকে সম্মানিত করেন।\lat{”}}

\textbf{আয়াত (সূরা আল-আ'রাফ ৭:৩২):}
\begin{quote}
\textbf{\lat{“}বলুন, আল্লাহ তাঁর বান্দাদের জন্য যে সৌন্দর্য ও পবিত্র রিজিক সৃষ্টি করেছেন, তা কে হারাম করেছে?\lat{”}} — তবে এই সৌন্দর্য অহংকারের জন্য নয়।
\end{quote}

\textbf{আয়াত (সূরা আল-কাহফ ১৮:৪৬):}
\begin{quote}
\textbf{\lat{“}ধন ও সন্তান দুনিয়ার জীবনের শোভা; কিন্তু স্থায়ী সৎকর্ম তোমার রবের কাছে প্রতিদান ও আশা হিসেবে শ্রেষ্ঠ।\lat{”}}
\end{quote}

\medskip\hrule\medskip

\subsection*{৭. গভীর শিক্ষা}

\textbf{১. \lat{power} (ক্ষমতা) বিনয় বাড়ায়।}
উমারের মতো সর্বোচ্চ ক্ষমতাবান মানুষ নিজেকে সবচেয়ে \lat{simple} (সাধারণ) রাখলেন — ক্ষমতা যত বাড়ে, \lat{humility} তত জরুরি।

\textbf{২. দুনিয়ার \lat{temptation} ত্যাগ।}
উমার রাজকীয় জীবন ছেড়ে \lat{simple} জীবন বেছে নিয়েছেন — দুনিয়ার চাকচিক্য ও প্রদর্শনের \lat{temptation} ত্যাগ করেছেন।

\textbf{৩. পোশাক — \lat{pride} বা \lat{humility}-র আয়না।}
পোশাক দিয়ে \lat{pride} করা বা মানুষকে নিচু দৃষ্টিতে দেখা — \lat{forbidden} (নিষিদ্ধ); \lat{humility}-র পোশাক মনকে নরম রাখে।

\textbf{৪. প্রদর্শনবর্জিত জীবন — মুমিনের \lat{character} (চরিত্র)।}
উমারের জীবন মানুষকে দেখানোর (\lat{display}) জন্য ছিল না; ছিল আল্লাহর জন্য। প্রদর্শনের (রিয়া) \lat{pride} মুমিনের নয়।

\textbf{৫. বায়তুল মাল ও আমানতদারির \lat{mission} (মিশন)।}
উমার রাষ্ট্রের কোষাগার রক্ষক ছিলেন; নিজের পোশাকের জন্য তার \lat{wealth} (সম্পদ) ব্যবহার করেননি — আমানতদারির চরম \lat{example}।

\medskip\hrule\medskip

\subsection*{৮. জীবনে প্রয়োগের পদ্ধতি}

\begin{enumerate}
\item \textbf{ক্ষমতায় বিনয়:} বড় পদ, ব্যবসা বা মর্যাদা পেলে সাধারণ জীবন বজায় রাখুন; মানুষের সাথে সহজ ও \lat{humble} (নম্র) থাকুন।
\item \textbf{পোশাকের দামে গর্ব নয়:} দামি পোশাক পরলেও তাতে \lat{pride} করবেন না; সস্তা পোশাকে লজ্জাবোধ করবেন না — উমার থেকে \lat{lesson} (শিক্ষা) নিন।
\item \textbf{প্রদর্শন \lat{avoid} (এড়িয়ে) চলুন:} সামাজিক মাধ্যমে বা বাস্তবে নিজের \lat{wealth}, পোশাক, ঘর-গাড়ি প্রদর্শন থেকে দূরে থাকুন।
\item \textbf{আমানতদারি:} রাষ্ট্র, প্রতিষ্ঠান বা পরিবারের সম্পদে নিজের জন্য বাড়তি নেবেন না — উমারের মতো \lat{just} (ন্যায়পরায়ণ) থাকুন।
\item \textbf{তাকওয়ার জন্য জীবন:} নিজের জীবনের লক্ষ্য ঠিক করুন — প্রদর্শনের \lat{pride} নয়, আল্লাহর নৈকট্য।
\end{enumerate}
\medskip\hrule\medskip

\subsection*{৯. রেফারেন্স}

\begin{itemize}
\item মুসান্নাফ ইবনে আবি শায়বাহ, হাদিস ৩৭১৯১ (সহীহ সনদ)
\item সহীহ মুসলিম, হাদিস ৯১ ও ৯১বি
\item সুনানে তিরমিজি, হাদিস ২৩৬৩
\item সূরা আল-আ'রাফ, আয়াত ৩২; সূরা আল-কাহফ, আয়াত ৪৬
\end{itemize}

\end{document}
